% Model: introduction of the layer principle, beside the layer-stack figure (scripts/plot_layers.py)
\section{Model}\label{sec:model}

\begin{figure}[h!]
  \centering
  \includegraphics[width=\textwidth]{coverage.pdf}
  \caption{Countries included in modelling and share of demand covered per level of modelling detail.
  Country borders from \citet{naturalearth}.}
  \label{fig:coverage}
\end{figure}

PyPSA-Earth \citep{pypsaearth} offers globally applicable data pipelines to compile grid, fleet, techno-economic assumptions and renewable potentials into runnable linear optimization energy models.
This project uses this pipeline for 132 countries, that are then run independently (Figure~\ref{fig:coverage}).
Sixteen of them are spatially resolved: the contiguous United States, Brazil, India, China and the twelve countries of North-West Europe whose extensive interconnection requires to be modelled within one model.
%
% [26]: the figure's height in text lines; without it wrapfig miscounts (parskip) and keeps narrowing lines onto the next page
\begin{wrapfigure}[26]{l}{0.47\textwidth}
  \vspace{-0.6\baselineskip}\centering
  \includegraphics[width=0.46\textwidth]{layers.pdf}
  \captionsetup{font=footnotesize}
  \caption{
  Data layers that go into each country's model (grid data is skipped for single-node countries) and the subsection structure this report contains.
}
  \label{fig:layers}
  \vspace{-0.8\baselineskip}
\end{wrapfigure}

Most countries share the same data pipeline to build the model, however local adjustments are needed where the default data source are insufficient.
In every subsection we describe both; the standard data processing, and individual adjustments on top.

In short, the PyPSA model finds the lowest-cost supply to meet demand exactly.
This formulation is accurate, if the data layers represent real power supply and demand.
We start with the grid which conditions and constrains power transport, and then set up the spatial and temporal distribution of demand.
Next, we build the supply side: the existing fleet, the spatial distribution of renewable potential, techno-economic assumptions and fuel costs.

